%! Parent Functions and Transformations — Unit 2, Lesson 2.2 — Guided Notes (Answer Key)
\documentclass[10pt]{article}
\usepackage{atda-article}
\usepackage{atda-key}   % pulls in -boxes; do NOT also load -boxes
\newcommand{\TallMath}[1]{$\displaystyle #1\rule[-1.4em]{0pt}{3.2em}$}
\setlength{\workrowsep}{6pt}

\begin{document}

\pageheader{Unit 2, Lesson 2.2}{Guided Notes --- Answer Key}
% NAMESTRIP: no \namedateperiod here — mirrors the blank. Teacher-only prose goes
% in the lesson plan as \begin{teachernote}[Guided Notes], never in a key.

\vspace{0.15in}

\begin{objectivebox}
\small
By the end of this lesson, I will be able to\dots
\begin{enumerate}[leftmargin=1.5em, itemsep=2pt, topsep=3pt]
  \item name the \textbf{parent function} of the linear, quadratic, and exponential families, and
        recognize each one from its equation, its graph, or a table of values;
  \item describe the transformation that turns a parent function into a given function, from the
        \textbf{equation} or from the \textbf{graph} --- translation, reflection, or dilation;
  \item explain why a number \emph{outside} the parentheses moves the graph the way you expect and
        a number \emph{inside} moves it the opposite way;
  \item say which of the \textbf{domain} and the \textbf{range} a transformation moved, and use
        that as a check on my answer.
\end{enumerate}
\end{objectivebox}

\vspace{0.10in}

\boxguard
\begin{vocabbox}
\small
Fill in each term as we build it together. Every one of them is a name for something you have
already done.
\vocabans{Parent function}{The simplest function in a family --- the one with no transformation
applied yet. Linear: $f(x)=x$. Quadratic: $g(x)=x^2$. Exponential: $h(x)=a^x$.}
\vocabans{Function family}{A parent function together with every transformation of it. All of them
share the parent's basic shape.}
\vocabans{Translation}{A slide. The graph moves up, down, left, or right; its shape and size do not
change. Written $f(x)+k$ (vertical) or $f(x+k)$ (horizontal).}
\vocabans{Reflection}{A flip across an axis. $-f(x)$ flips over the $x$-axis; $f(-x)$ flips over
the $y$-axis.}
\vocabans{Dilation}{A stretch or a squash. $k\,f(x)$ scales every output (vertical); $f(kx)$ scales
every input (horizontal).}
\end{vocabbox}

\vspace{0.10in}

\boxguard[12]
\begin{hookbox}
\small
\textbf{The same shell, three seconds later.} In Lesson 2.0 you read the graph of a firework
shell's height, $h_A(t)=-16t^2+64t$: it leaves the ground at $t=0$, tops out at $64$ feet when
$t=2$, and lands at $t=4$.

Tonight the crew fires \textbf{two identical shells}. The second one is launched \textbf{3 seconds
after} the first. Same shell, same climb, same fall --- just later. \textbf{Here is the question:
does the second shell's equation say $h_A(t+3)$ or $h_A(t-3)$?} Vote before you look at anything.

\ansline{$h_A(t-3)$. At $t=5$ the second shell has been up for $5-3=2$ seconds, so it is at}\\
\ansline{$h_A(2)=64$ feet. Later on the clock means \emph{subtract} inside the parentheses.}\\
\end{hookbox}

\vspace{0.10in}

\boxguard[30]
\begin{notesbox}{1. The Three Parent Functions}
\small
A \textbf{parent function} is the plainest member of a family --- the one with nothing done to it
yet. Every other function in the family is this graph, moved, flipped, or stretched.

\begin{center}
\begin{tabular}{@{}c@{\hspace{0.45cm}}c@{\hspace{0.45cm}}c@{}}
\begin{tikzpicture}
\begin{axis}[
    width=4.9cm, height=3.9cm, title={\scriptsize Linear},
    xmin=-3.4, xmax=3.4, ymin=-3.4, ymax=3.4,
    xtick={-3,0,3}, ytick={-3,0,3},
    grid=both, grid style={linegray!45}, axis lines=left,
    tick label style={font=\scriptsize}]
  \addplot[royal, very thick, domain=-3:3, samples=2]{x};
\end{axis}
\end{tikzpicture}
&
\begin{tikzpicture}
\begin{axis}[
    width=4.9cm, height=3.9cm, title={\scriptsize Quadratic},
    xmin=-3.4, xmax=3.4, ymin=0, ymax=9.6,
    xtick={-3,0,3}, ytick={0,3,6,9},
    grid=both, grid style={linegray!45}, axis lines=left,
    tick label style={font=\scriptsize}]
  \addplot[royal, very thick, domain=-3:3, samples=80]{x^2};
\end{axis}
\end{tikzpicture}
&
\begin{tikzpicture}
\begin{axis}[
    width=4.9cm, height=3.9cm, title={\scriptsize Exponential},
    xmin=-3.4, xmax=3.4, ymin=0, ymax=8.6,
    xtick={-3,0,3}, ytick={0,2,4,6,8},
    grid=both, grid style={linegray!45}, axis lines=left,
    tick label style={font=\scriptsize}]
  \addplot[royal, very thick, domain=-3:3, samples=80]{2^x};
\end{axis}
\end{tikzpicture}
\end{tabular}
\end{center}

\vspace{-4pt}
\renewcommand{\arraystretch}{1.5}
\begin{tabularx}{\linewidth}{@{}p{2.0cm} p{2.3cm} X p{2.5cm} p{2.5cm}@{}}
\rowcolor{mist}
\textbf{Family} & \textbf{Parent}\newline\textbf{equation} & \textbf{What its graph does} & \textbf{Domain}
    & \textbf{Range} \\
\hline
Linear & \ans{$f(x)=x$} & a straight line through the origin & \ans{$(-\infty,\infty)$}
    & \ans{$(-\infty,\infty)$} \\
Quadratic & \ans{$g(x)=x^2$} & a U-shaped curve (a parabola) with its turning point at the origin
    & \ans{$(-\infty,\infty)$} & \ans{$[0,\infty)$} \\
Exponential & \ans{$h(x)=2^x$} & hugs the horizontal axis on the left, then climbs faster and faster
    & \ans{$(-\infty,\infty)$} & \ans{$(0,\infty)$} \\
\end{tabularx}
{\scriptsize (Use $a=2$ for the exponential parent, as in the warm-up.)}

\vspace{6pt}
\textbf{Telling a line from an exponential when both climb forever} (Lesson 2.0's test): step the
input up by $1$ each time. A linear function \ans{adds} the same amount every step; an
exponential function \ans{multiplies} by the same amount every step.

\vspace{2pt}
And notice the exponential's range: $2^x$ gets closer and closer to $0$ but never reaches it, so
the range is $(0,\infty)$ --- a parenthesis, not a bracket. Lesson 2.5 gives that behaviour a name.
\end{notesbox}

\vspace{0.10in}

% [16] was tested here (to let the box split after the figure+table) and still gave
% 5 pages, so the break is the natural one. Restored to its box-sized value per 1.5.
\boxguard[30]
\begin{notesbox}{2. Sliding the Graph --- Translations}
\small
A \textbf{translation} slides a graph without changing its shape or size. Here is the quadratic
parent with two of its children.

\begin{center}
\begin{tikzpicture}
\begin{axis}[
    width=11.0cm, height=5.0cm,
    xlabel={$x$}, ylabel={$y$},
    xmin=-3.6, xmax=5.6, ymin=0, ymax=13,
    xtick={-3,-2,-1,0,1,2,3,4,5}, ytick={0,3,6,9,12},
    grid=both, grid style={linegray!45}, axis lines=left,
    tick label style={font=\scriptsize}, label style={font=\scriptsize},
    legend entries={{$y=x^2$},{$y=x^2+3$},{$y=(x-2)^2$}},
    legend style={font=\scriptsize, draw=none, fill=none,
      at={(0.5,1.02)}, anchor=south, legend columns=3}]
  \addplot[royal, very thick, domain=-3:3, samples=80]{x^2};
  \addplot[goldacc, very thick, dashed, domain=-3:3, samples=80]{x^2+3};
  \addplot[plumacc, very thick, dash dot, domain=-1:5, samples=80]{(x-2)^2};
\end{axis}
\end{tikzpicture}
\end{center}

\vspace{-2pt}
\renewcommand{\arraystretch}{1.5}
\begin{tabularx}{\linewidth}{@{}p{4.0cm} p{4.2cm} X@{}}
\rowcolor{mist}
\textbf{Written from the parent $f(x)=x^2$} & \textbf{Is the number outside or inside the
    parentheses?} & \textbf{Which way, and how far, the graph moves} \\
\hline
$g(x)=f(x)+3=x^2+3$ & \ans{outside} & \ans{up $3$} \\
$g(x)=f(x)-4=x^2-4$ & \ans{outside} & \ans{down $4$} \\
$g(x)=f(x-2)=(x-2)^2$ & \ans{inside} & \ans{right $2$} \\
$g(x)=f(x+5)=(x+5)^2$ & \ans{inside} & \ans{left $5$} \\
\end{tabularx}

\vspace{6pt}
\textbf{The rule, in two lines:}
\begin{itemize}[leftmargin=1.3em, itemsep=2pt, topsep=3pt]
  \item A number \ans{outside} the parentheses changes the \textbf{outputs}. The graph moves
        \textbf{up or down}, and it moves the way the sign tells you.
  \item A number \ans{inside} the parentheses changes the \textbf{inputs}. The graph moves
        \textbf{left or right}, and it moves the \ans{opposite} way from the sign you see.
\end{itemize}

\vspace{2pt}
\textbf{Why the opposite?} Read $f(x-2)$ at $x=5$: the rule first computes $5-2=3$, so the new
graph at $x=5$ shows whatever the old graph showed at $x=3$. Everything the parent did has been
copied $2$ units farther right.

\vspace{4pt}
\textbf{Back to the shells.} At $t=5$ the second shell has been in the air for $5-3=2$ seconds, so
$h_B(5)=h_A(2)=64$ feet. The rule is $h_B(t)=h_A(t-3)$ --- a \textbf{minus}, even though the shell
goes up \emph{later}. Its domain slid from $[0,4]$ to $[3,7]$ seconds; its range is still
$[0,64]$ feet.

\vspace{4pt}
\textbf{The cheapest check you have} (Lesson 2.1's words): a vertical slide moves the
\ans{range}; a horizontal slide moves the \ans{domain}.
\end{notesbox}

\vspace{0.10in}

\boxguard[30]
\begin{notesbox}{3. Flipping and Stretching --- Reflections and Dilations}
\small
A \textbf{reflection} flips a graph across an axis. A \textbf{dilation} stretches or squashes it.

\begin{center}
\begin{tabular}{@{}c@{\hspace{0.5cm}}c@{}}
\begin{tikzpicture}
\begin{axis}[
    width=7.4cm, height=4.6cm, title={\scriptsize Reflections of $h(x)=2^x$},
    xmin=-3.4, xmax=4.6, ymin=-9.5, ymax=9.5,
    xtick={-3,-2,-1,0,1,2,3}, ytick={-8,-4,0,4,8},
    grid=both, grid style={linegray!45}, axis lines=left,
    tick label style={font=\scriptsize}]
  \addplot[royal, very thick, domain=-3:3, samples=80]{2^x};
  \addplot[goldacc, very thick, dashed, domain=-3:3, samples=80]{-(2^x)};
  \addplot[plumacc, very thick, dash dot, domain=-3:3, samples=80]{2^(-x)};
  \node[font=\scriptsize, anchor=west, royal]  at (axis cs:3.1,8)   {$2^x$};
  \node[font=\scriptsize, anchor=west, goldacc] at (axis cs:3.1,-8)  {$-2^x$};
  \node[font=\scriptsize, anchor=west, plumacc] at (axis cs:3.1,0.9) {$2^{-x}$};
\end{axis}
\end{tikzpicture}
&
\begin{tikzpicture}
\begin{axis}[
    width=7.4cm, height=4.6cm, title={\scriptsize Dilations of $g(x)=x^2$},
    xmin=-3.4, xmax=4.6, ymin=0, ymax=13,
    xtick={-3,-2,-1,0,1,2,3}, ytick={0,3,6,9,12},
    grid=both, grid style={linegray!45}, axis lines=left,
    tick label style={font=\scriptsize}]
  \addplot[royal, very thick, domain=-3:3, samples=80]{x^2};
  \addplot[goldacc, very thick, dashed, domain=-2:2, samples=80]{3*x^2};
  \addplot[plumacc, very thick, dash dot, domain=-3:3, samples=80]{x^2/3};
  \node[font=\scriptsize, anchor=west, royal]  at (axis cs:3.1,9)  {$x^2$};
  \node[font=\scriptsize, anchor=west, goldacc] at (axis cs:2.1,12) {$3x^2$};
  \node[font=\scriptsize, anchor=west, plumacc] at (axis cs:3.1,3)  {$\tfrac{1}{3}x^2$};
\end{axis}
\end{tikzpicture}
\end{tabular}
\end{center}

\vspace{-4pt}
\renewcommand{\arraystretch}{1.5}
\begin{tabularx}{\linewidth}{@{}p{3.4cm} X@{}}
\rowcolor{mist}
\textbf{Notation} & \textbf{What it does to the graph} \\
\hline
$-f(x)$ & every output changes sign --- the graph flips over the \ans{$x$}-axis \\
$f(-x)$ & every input changes sign --- the graph flips over the \ans{$y$}-axis \\
$k \cdot f(x)$ with $|k|>1$ & every output is multiplied by $k$: a vertical \ans{stretch}
    (taller and steeper) \\
$k \cdot f(x)$ with $0<|k|<1$ & every output is multiplied by $k$: a vertical \ans{compression}
    (flatter) \\
$f(kx)$ with $|k|>1$ & a horizontal \emph{compression} --- the graph is squeezed toward the
    vertical axis \\
$f(kx)$ with $0<|k|<1$ & a horizontal \emph{stretch} --- the graph is pulled away from the
    vertical axis \\
\end{tabularx}

\vspace{6pt}
\textbf{Two things worth saying out loud.}
\begin{itemize}[leftmargin=1.3em, itemsep=2pt, topsep=3pt]
  \item $-f(x)$ and $f(-x)$ are \emph{different} flips. Warm-up item 3 already proved it:
        $-x^2$ gave $-9$ and $(-x)^2$ gave $9$.
  \item The parabola is its own reflection over the $y$-axis, because $(-x)^2=x^2$. A flip that
        leaves the picture unchanged is still a flip --- it just has nothing to show for it.
  \item $f(kx)$ is easiest to see when the domain is restricted. A machine that runs the same
        cycle in \emph{half} the time has the same graph squeezed onto $[0,4]$ instead of $[0,8]$.
\end{itemize}
\end{notesbox}

\vspace{0.10in}

\boxguard[24]
\begin{notesbox}{4. The Card: Which One Moved, the Domain or the Range?}
\small
Six notations, one question each: does the number touch the \textbf{outputs} or the
\textbf{inputs}? Fill the two blank columns and keep this page --- it is the whole lesson.

\renewcommand{\arraystretch}{1.5}
\begin{tabularx}{\linewidth}{@{}p{2.2cm} p{2.4cm} X p{2.7cm}@{}}
\rowcolor{mist}
\textbf{Notation} & \textbf{It changes the\dots} & \textbf{So the graph\dots}
    & \textbf{In context this moves the\dots} \\
\hline
$f(x)+k$ & \ans{outputs} & moves up or down, the way the sign says & \ans{range} \\
$f(x+k)$ & \ans{inputs} & moves left or right, the \emph{opposite} way & \ans{domain} \\
$-f(x)$ & \ans{outputs} & flips over the $x$-axis & \ans{range} \\
$f(-x)$ & \ans{inputs} & flips over the $y$-axis & \ans{domain} \\
$k \cdot f(x)$ & \ans{outputs} & stretches or compresses vertically & \ans{range} \\
$f(kx)$ & \ans{inputs} & stretches or compresses horizontally & \ans{domain} \\
\end{tabularx}

\vspace{6pt}
\textbf{Two errors to refuse before they happen:}
\begin{itemize}[leftmargin=1.3em, itemsep=2pt, topsep=3pt]
  \item ``$y=(x+5)^2$ moves right $5$.'' It moves \textbf{left} $5$. Inside the parentheses, the
        sign you see is the opposite of the direction you get.
  \item ``$-f(x)$ and $f(-x)$ both just flip it.'' They flip it over \emph{different} axes, and
        the warm-up showed they can give different numbers from the same input.
\end{itemize}
\end{notesbox}

\vspace{0.10in}

\boxguard[30]
\begin{practicebox}
\small
The robotics club sells raffle tickets. Selling $x$ hundred tickets gives a profit of
$P(x)=-x^2+8x$ hundred dollars, for $0 \le x \le 8$. Partway through, a sponsor promises a flat
\$400 gift no matter how many tickets sell, so the new profit is $Q$.

\begin{center}
\begin{tikzpicture}
\begin{axis}[
    width=10.8cm, height=4.8cm,
    xlabel={$x$ (tickets sold, hundreds)}, ylabel={profit (hundreds of dollars)},
    xmin=0, xmax=8.6, ymin=0, ymax=22,
    xtick={0,1,2,3,4,5,6,7,8}, ytick={0,4,8,12,16,20},
    grid=both, grid style={linegray!45}, axis lines=left,
    tick label style={font=\scriptsize}, label style={font=\scriptsize},
    legend entries={{$P(x)=-x^2+8x$},{$Q$}},
    legend style={font=\scriptsize, draw=none, fill=none,
      at={(0.5,1.02)}, anchor=south, legend columns=2}]
  \addplot[royal, very thick, domain=0:8, samples=80]{-x^2+8*x};
  \addplot[goldacc, very thick, dashed, domain=0:8, samples=80]{-x^2+8*x+4};
\end{axis}
\end{tikzpicture}
\end{center}

\begin{enumerate}[leftmargin=1.6em, itemsep=6pt, topsep=2pt]
  \item Describe how the graph of $Q$ sits compared with the graph of $P$, then write $Q$ in
        function notation, starting from $P$.

\ansline{Same shape, shifted up $4$ --- a vertical translation of $4$ hundred dollars.}\\
\ansline{$Q(x)=P(x)+4$, that is $Q(x)=-x^2+8x+4$.}\\

  \item State the domain and the range of $P$, then of $Q$, in interval notation. Which one moved
        --- the domain or the range --- and why is that exactly what you should have expected?

\ansline{$P$: domain $[0,8]$, range $[0,16]$. \quad $Q$: domain $[0,8]$, range $[4,20]$.}\\
\ansline{The \textbf{range} moved. The $+4$ sits outside the parentheses, so it changes}\\
\ansline{outputs; how many tickets they may sell did not change at all.}\\

  \item Describe the transformation of the \emph{parent} function each of these is. Name the
        parent first. \quad (a) $y=x^2-6$ \quad (b) $y=(x+4)^2$ \quad (c) $y=-2^x$

\ansline{(a) parent $x^2$; translated down $6$.}\\
\ansline{(b) parent $x^2$; translated left $4$ (inside the parentheses, opposite sign).}\\
\ansline{(c) parent $2^x$; reflected over the $x$-axis.}\\

  \item Check your reading of the highest point against the arithmetic: find $P(4)$, then $Q(4)$.
\begin{work}
  P(4) &= -(4)^2+8(4) \\
       &= -16+32 \\
       &= 16 \\
  Q(4) &= P(4)+4 \\
       &= 20
\end{work}
\end{enumerate}
\end{practicebox}

\end{document}
